Replace the heights by ranks $1,\ldots,5$. The adjacent comparisons must alternate. The positions of $1$ and $5$ differ by $3$, so they occupy $(1,4)$ or $(2,5)$, in either order.

If $1$ is first and $5$ fourth, the arrangement is $(1,a,b,5,c)$, where $(a,b,c)$ permutes $(2,3,4)$. Alternation requires only $a>b$, giving $3$ choices. Reversal and height complementation send this case bijectively to the other three placements. Thus the answer is $\boxed{4\cdot3=12}$.
